% from ega1-9-fr.tex
\oldpage[I]{175}
il en résulte que $\mathscr{F}$ est limite inductive de ses
sous-$\mathscr{O}_X$-Modules quasi-cohérents de type fini. En effet, pour
tout ouvert affine $U\subset X$, on a $\mathscr{F}|U=\widetilde{M}$, où $M$
est un $A(U)$-module, et comme ce dernier est limite inductive de ses
sous-modules de type fini, $\mathscr{F}|U$ est limite inductive de ses
sous-$(\mathscr{O}_X|U)$-Modules quasi-cohérents de type fini
(\hyperref[I.1.3.9-fr]{1.3.9}). Or, par hypothèse, chacun de ces sous-Modules
est induit sur $U$ par un sous-$\mathscr{O}_X$-Module quasi-cohérent de type
fini $\mathscr{G}_{\lambda,U}$ de $\mathscr{F}$. Les sommes finies des
$\mathscr{G}_{\lambda,U}$ sont encore des $\mathscr{O}_X$-Modules
quasi-cohérents de type fini, car la question est locale et le cas où $X$
est affine a été traité dans (\hyperref[I.1.3.10-fr]{1.3.10})~; il est clair
alors que $\mathscr{F}$ est limite inductive de ces sommes finies, d'où notre
assertion.
\end{remark}

\begin{corollary}[9.4.5]
\label{I.9.4.5-fr}
Sous les hypothèses de (9.4.3), pour tout $(\mathscr{O}_X|U)$-Module
quasi-cohérent de type fini $\mathscr{G}$, il existe un
$\mathscr{O}_X$-Module quasi-cohérent de type fini $\mathscr{G}'$ tel que
$\mathscr{G}'|U=\mathscr{G}$.
\end{corollary}

\begin{proof}
Comme $\mathscr{F}=j_*(\mathscr{G})$ est quasi-cohérent (9.4.2) et
$\mathscr{F}|U=\mathscr{G}$, il suffit d'appliquer (9.4.3) à $\mathscr{F}$.
\end{proof}

\begin{lemma}[9.4.6]
\label{I.9.4.6-fr}
Soient $X$ un préschéma, $L$ un ensemble bien ordonné,
$(V_\lambda)_{\lambda\in L}$ un recouvrement de $X$ par des ouverts affines,
$U$ un ouvert de $X$~; pour tout $\lambda\in L$, on pose
$W_\lambda=\bigcup_{\mu<\lambda}V_\mu$. On suppose que~:
\begin{enumerate}
  \item[1\textsuperscript{o}] Pour tout $\lambda\in L$,
    $V_\lambda\cap W_\lambda$ soit quasi-compact~;
  \item[2\textsuperscript{o}] Le morphisme d'immersion $U\to X$ est
    quasi-compact.
\end{enumerate}
Alors, pour tout $\mathscr{O}_X$-Module quasi-cohérent $\mathscr{F}$ et tout
sous-$(\mathscr{O}_X|U)$-Module quasi-cohérent de type fini $\mathscr{G}$ de
$\mathscr{F}|U$, il existe un sous-$\mathscr{O}_X$-Module quasi-cohérent de
type fini $\mathscr{G}'$ de $\mathscr{F}$ tel que
$\mathscr{G}'|U=\mathscr{G}$.
\end{lemma}

\begin{proof}
Posons $U_\lambda=U\cup W_\lambda$~; on va définir par récurrence transfinie
une famille $(\mathscr{G}'_\lambda)$, où $\mathscr{G}'_\lambda$ est un
sous-$(\mathscr{O}_X|U_\lambda)$-Module quasi-cohérent de type fini de
$\mathscr{F}|U_\lambda$, tel que
$\mathscr{G}'_\lambda|U_\mu=\mathscr{G}'_\mu$ pour $\mu<\lambda$ et
$\mathscr{G}'_\lambda|U=\mathscr{G}$. L'unique sous-$\mathscr{O}_X$-Module
$\mathscr{G}'$ de $\mathscr{F}$ tel que
$\mathscr{G}'|U_\lambda=\mathscr{G}'_\lambda$ pour tout $\lambda\in L$
(\hyperref[0.3.3.1-fr]{0, 3.3.1}) répondra à la question. Supposons donc les
$\mathscr{G}'_\mu$ définis et ayant les propriétés précédentes pour
$\mu<\lambda$~; si $\lambda$ n'a pas de prédécesseur on prendra pour
$\mathscr{G}'_\lambda$ l'unique sous-$(\mathscr{O}_X|U_\lambda)$-Module de
$\mathscr{F}|U_\lambda$ tel que
$\mathscr{G}'_\lambda|U_\mu=\mathscr{G}'_\mu$ pour tout $\mu<\lambda$, ce
qui est licite puisque les $U_\mu$ avec $\mu<\lambda$ forment alors un
recouvrement de $U_\lambda$. Si au contraire $\lambda=\mu+1$, on a
$U_\lambda=U_\mu\cup V_\mu$, et il suffira de définir un
sous-$(\mathscr{O}_X|V_\mu)$-Module quasi-cohérent de type fini
$\mathscr{G}''_\mu$ de $\mathscr{F}|V_\mu$ tel que
\[
  \mathscr{G}''_\mu|(U_\mu\cap V_\mu)
  =\mathscr{G}'_\mu|(U_\mu\cap V_\mu)~;
\]
on prendra ensuite pour $\mathscr{G}'_\lambda$ le
sous-$(\mathscr{O}_X|U_\lambda)$-Module de $\mathscr{F}|U_\lambda$ tel que
$\mathscr{G}'_\lambda|U_\mu=\mathscr{G}'_\mu$ et
$\mathscr{G}'_\lambda|V_\mu=\mathscr{G}''_\mu$
(\hyperref[0.3.3.1-fr]{0, 3.3.1}). Or, comme $V_\mu$ est affine, l'existence
de $\mathscr{G}''_\mu$ est assurée par (9.4.3) dès que l'on a démontré que
$U_\mu\cap V_\mu$ est quasi-compact~; mais $U_\mu\cap V_\mu$ est réunion de
$U\cap V_\mu$ et de $W_\mu\cap V_\mu$, qui sont tous deux quasi-compacts en
vertu de l'hypothèse.
\end{proof}

\begin{theorem}[9.4.7]
\label{I.9.4.7-fr}
Soient $X$ un préschéma, $U$ un ouvert de $X$. On suppose vérifiée l'une des
conditions suivantes~:
\begin{enumerate}
  \item[(a)] L'espace sous-jacent à $X$ est localement noethérien.
  \item[(b)] $X$ est un schéma quasi-compact et $U$ un ouvert quasi-compact.
\end{enumerate}
Pour tout $\mathscr{O}_X$-Module quasi-cohérent $\mathscr{F}$ et tout
sous-$(\mathscr{O}_X|U)$-Module quasi-cohérent de type fini $\mathscr{G}$ de
$\mathscr{F}|U$, il existe alors un sous-$\mathscr{O}_X$-Module
quasi-cohérent de type fini $\mathscr{G}'$ de $\mathscr{F}$ tel que
$\mathscr{G}'|U=\mathscr{G}$.
\end{theorem}

